\documentclass[10pt,a4paper]{amsart}
\usepackage[margin=19mm]{geometry}
\usepackage{amsmath,amssymb,mathtools}
\usepackage[scr=rsfs,cal=euler]{mathalfa}
\usepackage{xfrac}
\usepackage[unicode,hidelinks,bookmarks=false]{hyperref}
\newcommand{\ie}{i.e.\ }
\newcommand{\cf}{cf.\ }
\newcommand{\resp}{respectively\ }
\newcommand{\Subsection}{\subsection}
\providecommand{\nobreakdash}{\nobreak\hbox{-}}
\setlength{\emergencystretch}{2em}
\multlinegap0pt
\usepackage{float}
\usepackage{csquotes}
\usepackage{tikz}
\newtheorem{lemma}{Lemma}
\newtheorem{remark}{Remark}
\newtheorem{proposition}{Proposition}
\newcommand{\apyr}[1][l]{\ensuremath{pyramid(#1)}}
\newcommand{\chin}[2][n]{\ensuremath{ch[#1;#2]}}
\newcommand{\pyr}[3][l]{\ensuremath{pyramid(#1,#2,#3)}}
\title{Polychrony as Chinampas}
\author{Eric Dolores-Cuenca, Jose Antonio Arciniega-Nevarez, Anh Nguyen, Yitong Zou, Luke Van Popering, Nathan Crock, Gordon Erlebacher, Jose L. Mendoza-Cortes}
\date{Source: arXiv:2103.15265v4 (CC BY 4.0)}
\begin{document}
\maketitle

\noindent\emph{Source and licence.} Polychrony as Chinampas. Version arXiv:2103.15265v4 (CC BY 4.0), distributed by arXiv under the Creative Commons Attribution 4.0 International licence, http://creativecommons.org/licenses/by/4.0/, as recorded in the arXiv licence metadata for this exact version and shown on the abstract page https://arxiv.org/abs/2103.15265v4. Full licence notice: \texttt{licenses/arxiv-2103.15265-CC-BY-4.0.txt}.\par\smallskip

\noindent\textit{Editorial scope.} Counted unit: the article's proposition lemma:one (source0.tex lines 581--587). The article's complete original proof of it is delivered with it (source0.tex lines 589--630, one \texttt{proof} environment of 42 lines). No mathematics is written, repaired or re-derived: the statement and the proof are the article's own lines, in the article's own notation, and no retained line is substituted or re-flowed.  Retained uncounted as the source-local context the proof needs: lines 473--475; lines 488--490; lines 539--541; lines 543--550.  One declared adaptation: exactly 1 ancillary strings inside the retained spans are omitted (image-inclusion commands and, where the source repeats a label, the repeated label definitions) and no other line differs from the source; the figure environments, with their placement, captions and labels, are kept, and the package records the omitted strings in fidelity receipts bound to the affected bodies.  No retained line contains a citation, so the wrapper carries no bibliography and no bibliography entry.  The retained lines contain the article's own diagram code, which the wrapper typesets with the standard tikz packages; no image file is delivered and the package declares no asset dependency.  Editorial formatting only, in addition: an AMS \texttt{amsart} preamble that restates the article's own declarations and macros used by the retained lines, namely the environments \texttt{lemma}, \texttt{remark}, \texttt{proposition} on separate counters, together with the standard packages those lines need.  The retained environments print in this document as Lemma 1, Remark 1, Proposition 1 in order of appearance, whereas the article prints the same items with the numbering of its own full document.  The complete native source of the article as distributed in the arXiv e-print for this exact version is bundled unchanged as \texttt{sources/109658/source0.tex}.
\begin{lemma}\label{lemma:plus}
Let $C$ be a chinampa with profit $n$. Stacking a $\pyr{i}{t}$ with $l\geq 3$ into $C$ creates a chinampa $C'$ with profit $m>n$.  
\end{lemma}

\begin{remark}\label{rmk:prof-0-1}
Stacking one pyramid with a length of three and several pyramids with a length of two returns a chinampa with a profit of zero. Conversely, under Lemma \ref{lemma:plus}, if we have a chinampa with a profit of zero, then it must be the result of stacking one pyramid with a length of three and some pyramids with a length of two. Similarly, a chinampa with a profit of one has  two copies of pyramids with a length of three and several copies of pyramids with a length of two. For a profit of two,  we can have either one $\apyr[4]$ and several $\apyr[2]$ instances stacked on or below it or three  $\apyr[3]$ instances and several $\apyr[2]$ instances.
\end{remark}

\begin{equation}
p(x)= \frac{d}{dx} \left(e^{2x}\int f(x) +h(x)e^{-2x} dx\right)\label{eq:sol}.
\end{equation}

\begin{proposition}
\label{lemma:zero} In $\apyr[n+3]$, for the nonnegative integer $n$, there are $(2+3n) 2^{n-1}$ possible zero-profit chinampas. Furthermore, these numbers are given by the coefficients of the generating function


\begin{equation*}
\sum_{n=0}^\infty \chin[n+3]{(3,1)}\frac{x^n}{n!}=(3x+1)e^{2x}.    
\end{equation*}
\end{proposition}

\begin{proposition}\label{lemma:one}
Chinampas of a certain unit of profit have the generating function

\begin{equation*}
\sum_{n=0}^{\infty} \chin[n+4]{ (3,2)}\frac{x^n}{n!}=(9x^2+18x+4)\frac{e^{2x}}{2}.
\end{equation*}
\end{proposition}

\begin{proof} 
A chinampa unit of profit can only be formed by two copies of $\apyr[3]$ and chains of $\apyr[2]$, as shown in Remark~\ref{rmk:prof-0-1}. We define the formal series

 \[q(x) = \sum_{n=0}^{\infty} \chin[n+4]{(3,2)}\frac{x^n}{n!},\] 
and consider the following cases:

\begin{itemize}
\item None of the instances of $\apyr[3]$ are at the top. We then have subpyramids as in the previous proposition, so we count $2\chin[n+3]{(3,2)}.$ 
\item One $\apyr[3]$ instance is at the top. Then, the remaining $\apyr[3]$ instances can be placed in 
    $2\chin[n+3]{(3,1)}$ ways on the two subpyramids. However, the two subcases have $\chin[n+2]{(3,1)}$ terms in common, as shown in Figure \ref{fig:double}. Thus, the correct number of combinations is $2\chin[n+3]{ (3,1)}-\chin[n+2]{(3,1)}$.
\end{itemize}

\begin{figure}[H]

\caption{Stacking any pyramid onto pyramid(3) in the vertex, where the red and blue lines collide and continue the process of stacking pyramids iteratively. This process describes a family of pyramids counted twice: once under the blue region and once under the red region. }
\label{fig:double}
\end{figure}

We conclude that for each nonnegative integer $n$, we have

\[\chin[n+4]{(3,2)}=2\chin[n+3] {(3,2)}+2\chin[n+3]{ (3,1)}-\chin[n+2]{(3,1)}.\]

Then, using the generating series $p(x)$ found in Lemma~\ref{lemma:zero}, we compute

\begin{eqnarray*}
q(x)&=&
\sum_{n=0}^{\infty} \chin[n+4]{ (3,2)}\frac{x^n}{n!}\\
&=&2\sum_{n=1}^{\infty} \chin[n+3] {(3,2)}\frac{x^n}{n!}+2\sum_{n=0}^{\infty} \chin[n+3] {(3,1)}\frac{x^n}{n!}-\sum_{n=1}^{\infty} \chin[n+2]{(3,1)}\frac{x^n}{n!}\\
&=&2\int q(x)dx+2 (3x+1)e^{2x}- \int (3x+1)e^{2x}dx\\
&=&2\int q(x)dx+ \frac{18x+9}{4}e^{2x}
\end{eqnarray*}

By solving Equation (\ref{eq:sol}), we obtain

\begin{eqnarray*}
\sum_{n=0}^{\infty} \chin[n+4] {(3,2)}\frac{x^n}{n!}&=&\frac{9+2c+36x+18x^2}{4}e^{2x}
\end{eqnarray*}

We compute $\chin[4]{(3,2)}=2$ to conclude

\[\sum_{n=0}^{\infty} \chin[n+4] {(3,2)}\frac{x^n}{n!}=(4+18x+9x^2)\frac{e^{2x}}{2}.\]
\end{proof}

\end{document}
